\documentclass{styles/jouthesis}

\JOUTitleCN{一种支持并发执行的脚本编程语言设计与编译器实现}
\JOUTitleEN{A Tcl-style CPU/CUDA Hybrid Computational Geometry Scripting System}
\JOUSetAuthor{房鸿铭}
\JOUSetStudentID{2022122407}
\JOUSetCollege{计算机学院}
\JOUSetMajor{软件工程}
\JOUSetAdvisor{纪兆辉}
\JOUSetDate{May 2026}

\JOUSharedApplyToThesis

\begin{document}

\frontmatter
\makecover

\begin{cnabstract}
    许多算法原型止步于两个断点：脚本表达足够灵活，却难以触达 GPU；底层 CUDA 足够强大，却把使用者推向繁琐的内存管理和调度细节。本文围绕这一矛盾设计并实现 FTCL，一个 Tcl 风格的算法实验脚本系统，并将其扩展为面向计算几何的 CPU/CUDA 混合执行平台。FTCL 以 C++20 实现，支持变量、过程、作用域、列表、字典、数组、表达式、命令替换、源码加载、测试命令和线程通道。它的目标不是复刻完整 Tcl，而是提供一个足够小、足够清晰、又足以承载真实算法实验的控制层，让复杂计算可以被脚本组织、复现和扩展。

    系统核心是跨设备向量抽象 UVec。UVec 将一个逻辑向量映射为 CPU 缓冲区和多个 CUDA 缓冲区，通过有效标记维护最新副本，并在 CPU 或 CUDA 请求数据时自动完成同步。它还提供统一原始指针接口，使 CPU 算法和 CUDA kernel 能在同一命令模型下工作。基于 UVec，本文实现 geom 命令库，覆盖二维几何基础操作以及距离矩阵、最近点、圆形范围计数、包围盒归约、仿射变换和 AABB 碰撞检测等批量算法。实现中进一步将 UVec 句柄管理从全局锁细化为按句柄加锁，降低异构并发中的不必要串行化。

    本文从正确性、性能和应用三个层面评估 FTCL。语义回归、CPU/GPU 等价、随机差分和边界测试验证系统可靠性；单 GPU 与八卡 A800 多 GPU 实验显示，紧凑输出的范围计数可接近线性弱扩展，而距离矩阵受输出分配和内存带宽限制；并发实验揭示解释器调度、锁粒度和 GPU 资源竞争对吞吐量与尾延迟的影响。最后，本文将点、圆、包围盒和距离矩阵映射到海洋牧场感知、智慧港口安全、无人机水面巡检和船舶碰撞检测等场景。实验表明，FTCL 不只是几何函数集合，而是一个把脚本表达、跨设备内存、并发执行和可复现实验连接起来的异构空间计算原型。

\cnkeywords{脚本解释器；Tcl；计算几何；UVec；CUDA；多 GPU；海洋空间计算}
\end{cnabstract}

\begin{enabstract}
    Algorithm prototypes often break at the boundary between expression and execution: scripts are easy to write but rarely reach the GPU, while CUDA is powerful but forces users to manage memory, pointers, and synchronization by hand. This thesis addresses that gap by designing FTCL, a Tcl-style scripting system for algorithm experiments, and extending it into a CPU/CUDA hybrid computational geometry platform. Implemented in C++20, FTCL supports variables, procedures, scopes, lists, dictionaries, arrays, expressions, command substitution, source loading, testing commands, and thread channels. Its purpose is not to reproduce the whole Tcl language, but to provide a compact control layer that can organize, reproduce, and extend real computational workloads.

    The central abstraction is UVec, a cross-device vector that represents one logical vector through CPU and CUDA buffers. UVec tracks which copy is valid, synchronizes data on demand when a device requests the latest value, and exposes unified raw pointer interfaces for CPU algorithms and CUDA kernels. On top of UVec, the thesis implements a script-facing \texttt{geom} command library, including basic two-dimensional geometry operations and batch algorithms such as distance matrices, nearest point queries, circle range counting, bounding-box reduction, affine transformations, and AABB collision checks. The runtime also refines UVec handle management from one global lock to per-handle locking, reducing unnecessary serialization in heterogeneous concurrent workloads.

    The system is evaluated from correctness, performance, and application perspectives. Semantic regression tests, CPU/GPU equivalence tests, random differential tests, and edge-case tests validate reliability. Single-GPU and eight-GPU A800 experiments show that compact-output workloads such as circle range counting can achieve near-linear weak scaling, while output-heavy workloads such as distance matrices are limited by allocation and memory bandwidth. Concurrency experiments further reveal how interpreter scheduling, lock granularity, and GPU contention affect throughput and tail latency. Finally, the thesis maps points, circles, bounding boxes, and distance matrices to ocean-oriented scenarios such as marine ranch sensing, smart port safety, UAV water-surface inspection, and ship collision checking. The result is not merely a geometry library, but a prototype heterogeneous spatial-computing system that connects scripting, cross-device memory, concurrent execution, and reproducible evaluation.

\enkeywords{Scripting interpreter; Tcl; Computational geometry; UVec; CUDA; Multi-GPU; Marine spatial computing}
\end{enabstract}

\tableofcontents
\JOUListOfTablesAndFigures

\mainmatter
\chapter{Introduction}

\section{Background}
Scripting languages are often used as glue layers between low-level system capabilities and high-level application logic. Tcl is a representative example of this design philosophy. Its core syntax is simple: a script is a sequence of commands, and each command is a sequence of words. Variable substitution, command substitution, and expression evaluation are performed at runtime \cite{ousterhout1994tcl,tcldocs}. This command-centered model is easy to embed and easy to extend with new native commands.

Computational geometry is another important area for algorithm-oriented software. Problems involving points, segments, polygons, bounding boxes, distances, intersections, nearest-neighbor queries, and range queries appear in programming contests, graphics, robotics, games, and map analysis \cite{deberg2008computational,preparata1985computational}. Many geometry workloads are also naturally parallel. For example, each entry of a distance matrix can be computed independently, and each range query can be assigned to a different thread. This makes computational geometry a suitable domain for GPU acceleration \cite{cuda-guide}.

This thesis combines these two ideas and places them in an ocean-oriented application context. FTCL uses a Tcl-style interpreter as the control layer, C++ as the host implementation language, and CUDA as the acceleration backend for selected batch geometry algorithms. The goal is not merely to implement geometry functions, but to build a scriptable system where a user can describe marine ranch sensing, smart port spatial analysis, UAV inspection, or ship collision checking without manually managing CUDA memory, raw device pointers, or synchronization.

\section{Motivation}
The project is motivated by three observations.

First, an interpreter project is a good way to demonstrate systems engineering. It requires parsing, value representation, scoping, command dispatch, error handling, testing, and performance evaluation. These components interact heavily, so the project is richer than a single standalone algorithm implementation.

Second, geometry algorithms are familiar to programming-contest students and easy to validate visually. Basic functions such as orientation and segment intersection are small, while batch tasks such as distance matrices and nearest-point queries are large enough to show performance scaling.

Third, heterogeneous computing needs a practical data abstraction. If every geometry command manually calls cudaMalloc, cudaMemcpy, and cudaFree, the script layer becomes difficult to use and difficult to test. UVec addresses this problem by providing a vector-like object that can hold both CPU and CUDA copies and synchronize them on demand.

Fourth, many realistic spatial-computing tasks are no longer single-device tasks. A marine monitoring or port inspection system may involve many sensors, large detection sets, and multiple independent query regions. If the hardware platform provides several GPUs, the scripting system should be able to expose these devices and dispatch independent workloads without being serialized by the interpreter runtime.

\section{Research Questions}
This thesis is organized around the following research questions.

\begin{description}
  \item[RQ1:] Can a compact Tcl-style interpreter provide a stable enough semantic foundation for algorithmic scripts, concurrency primitives, and extension commands?
  \item[RQ2:] Can a cross-device vector abstraction hide CPU/CUDA data movement while preserving a simple vector-like programming model?
  \item[RQ3:] Which computational geometry workloads benefit from CUDA acceleration when the measurement is end-to-end script command latency rather than pure kernel time?
  \item[RQ4:] What runtime overheads limit the acceleration of a script-driven heterogeneous geometry system, and how can these overheads be observed experimentally?
  \item[RQ5:] Can the same scripting and UVec model express multi-GPU spatial workloads, and what limits the scaling of different geometry algorithms?
\end{description}

These questions separate the work into five layers: language semantics, memory abstraction, algorithm acceleration, multi-GPU execution, and experimental analysis. This distinction is important because a successful scripting system must be correct and usable before raw performance becomes meaningful.

\section{Main Contributions}
The main contributions of this thesis are as follows.

\begin{enumerate}
  \item \textbf{A script runtime for algorithm experiments.} FTCL implements a Tcl-style command model, variable substitution, expression evaluation, procedures, collections, source loading, testing commands, and thread channels in C++20.
  \item \textbf{A cross-device vector abstraction.} UVec models one logical vector with CPU and CUDA physical copies. It tracks valid copies, performs synchronization on demand, and exposes unified read-only and mutable raw pointer interfaces.
  \item \textbf{A finer-grained UVec handle manager.} The runtime avoids executing long geometry commands while holding a global UVec table lock. The global mutex protects handle-table lookup, while per-handle entry locks protect individual UVec objects. This design allows independent UVec workloads to run concurrently on different CUDA devices.
  \item \textbf{A script-facing computational geometry library.} The \texttt{geom} command family exposes scalar and batch geometry operations to FTCL scripts, allowing large point sets and query sets to be passed by UVec handles instead of repeated script strings.
  \item \textbf{CUDA acceleration for batch geometry.} The project implements CUDA backends for representative independent-work-item algorithms such as distance matrix, nearest point, range counting, transformation, orientation, and AABB collision.
  \item \textbf{A multi-GPU scaling experiment.} The evaluation uses an eight-GPU A800 server to measure 1/2/4/8 GPU weak scaling. The result shows near-linear scaling for compact-output range counting, partial scaling for nearest-point search, and output-bandwidth limitations for distance matrices.
  \item \textbf{An ocean-oriented application framing.} The thesis maps points, circles, boxes, and distance matrices to marine ranch sensing, smart port spatial computing, UAV floating-object inspection, and ship-path collision checking.
  \item \textbf{A reproducible evaluation pipeline.} The project includes semantic tests, CPU/GPU equivalence tests, random differential tests, edge-case tests, performance-scale tests, ablation studies, concurrency experiments, multi-GPU CSV outputs, and generated thesis figures.
\end{enumerate}

The novelty of the project lies in the integration of these components. FTCL is not only an interpreter, UVec is not only a container, and the geometry library is not only a set of algorithms. Together they form a script-driven heterogeneous geometry experimentation platform.

\section{Thesis Organization}
Chapter 2 introduces related technologies. Chapter 3 presents the overall system design. Chapter 4 describes the implementation of the interpreter, UVec, and geometry backends. Chapter 5 evaluates correctness and performance. Chapter 6 discusses applications, limitations, and possible improvements. The appendix gives build, test, benchmark, and thesis compilation commands.

\chapter{Related Work and Technical Background}

\section{Tcl-style Interpretation}
A Tcl-style language is command-oriented. The first word of a command is used as the command name, and the remaining words are passed as arguments after substitution. This model differs from expression-centered languages. In a Tcl-like system, most syntax is deliberately shallow, while semantics are provided by commands.

FTCL follows this philosophy but implements only the subset needed for the project. It supports variables, arrays, procedures, control flow, lists, dictionaries, expression evaluation, source loading, and tests. The system is not intended to be a complete Tcl implementation. Instead, it is a controlled interpreter used as a platform for geometry and heterogeneous-computing experiments.

\section{Parser Design}
FTCL currently contains two parser directions. The legacy parser reads directly from a character stream. It has a small constant factor and is useful as a stable reference. The token-stream parser first performs lexing and then uses token cursors to identify command and word boundaries. The token-stream path is not always faster, but it is more suitable for future work such as span-based diagnostics, AST construction, bytecode generation, and parser backend comparison.

The project also uses shadow comparison during migration. In this mode, both parser paths can parse the same input and compare the resulting structure. This reduces the risk of semantic regressions when replacing old parsing logic with a more structured token-stream implementation.

\section{Expression Evaluation}
Expression evaluation is handled separately from ordinary command parsing. The \texttt{expr} command needs operator precedence, associativity, short-circuit logic, ternary expressions, integer overflow checks, bitwise operations, division semantics, variable references, and command substitution. FTCL uses layered recursive descent functions, where lower-precedence functions call higher-precedence functions. This structure makes expressions such as \texttt{a || b \&\& c} parse correctly because logical AND is handled below logical OR.

Short-circuit behavior is important. When the left side of a logical OR is true, the right side must still be parsed so that the parser position remains correct, but it should not execute side effects. FTCL handles this using an evaluation mode guard.

\section{Computational Geometry}
Computational geometry studies algorithms for geometric objects such as points, segments, polygons, and regions. This project focuses on two-dimensional geometry. The basic operations include dot product, cross product, orientation, distance, segment intersection, point-to-segment distance, polygon area, point-in-polygon classification, and convex hull construction.

Batch geometry algorithms are particularly relevant to GPU acceleration. A distance matrix has one output per point pair. A range-count query has one output per query region. A nearest-point query has one output per query point. These operations contain many independent work items and can therefore be mapped naturally to CUDA threads.

\section{CUDA and End-to-end Performance}
CUDA exposes the GPU as a massively parallel device. A typical CUDA program allocates device memory, transfers input data, launches kernels, synchronizes, and copies results back to the host. However, a user of a scripting language does not want to manage these steps manually. FTCL therefore hides them behind UVec and geom commands.

The experiments in this thesis measure end-to-end FTCL command time rather than pure kernel time. This includes command dispatch, handle lookup, UVec synchronization, output allocation, kernel launch, synchronization, and sometimes readback. End-to-end timing is more conservative, but it better reflects what a script user experiences.

\section{Reproducible Experiment Artifacts}
The project stores benchmark scripts, CSV files, and SVG/PNG figures under the \texttt{docs} directory. This organization is important for a thesis project because it links claims to reproducible artifacts. Instead of manually drawing figures, the project regenerates them from executable tests and benchmark data.

\chapter{系统设计}

\section{总体架构}
FTCL 的总体架构如图 \ref{fig:system-arch} 所示。一个脚本首先被解析为 commands 和 words，随后由 interpreter 对 word 进行求值，完成变量替换、命令替换、数组访问和字符串拼接，最后将第一个求值结果作为 command name 进行分发。\texttt{geom} command 位于语言层和计算层之间：它负责解析几何参数或 UVec handle，并根据用户指定的 device 选择 CPU backend 或 CUDA backend。UVec 则处在 command layer 和 backend layer 之间，负责维护跨设备数据的位置、有效性和同步过程。

\begin{figure}[H]
  \centering
  \includegraphics[width=0.96\textwidth]{system-architecture.png}
  \caption{FTCL 几何脚本系统总体架构（FTCL geometry scripting system architecture）}
  \label{fig:system-arch}
\end{figure}

这个设计的核心是分层而不是堆叠。Parser 和 interpreter 专注于语言语义；command layer 提供稳定的用户接口；UVec 负责 data residency 和 synchronization；CPU/CUDA backend 只实现具体算法。这样做的好处是：语言可以继续扩展，几何算法可以继续增加，CUDA 路径也可以在不破坏脚本接口的情况下逐步优化。

\section{工程规模}
表 \ref{tab:engineering-scope} 总结了项目的主要工程组成。列出这张表的目的不是简单展示文件数量，而是让系统边界变得清楚：FTCL 不是几个孤立的几何函数，而是由 interpreter、memory abstraction、geometry library、CUDA backend、testing framework 和 benchmark pipeline 共同组成的完整原型系统。

\begin{table}[htbp]
\centering
\caption{FTCL 系统工程组成（Engineering scope）}
\label{tab:engineering-scope}
\small
\begin{tabularx}{\textwidth}{lXXX}
\toprule
Component & Main files & Main classes/functions & Role \\
\midrule
Interpreter core & \texttt{interp.hpp}, \texttt{type.hpp}, \texttt{parser.hpp} & \texttt{Interp}, \texttt{Value}, \texttt{Word}, \texttt{Script} & 负责 script parsing、evaluation、variables、procedures 和 command dispatch。 \\
Expression parser & \texttt{expr.hpp} & \texttt{ExprParser}, \texttt{Datum} & 支持 arithmetic、comparison、logical、bitwise、ternary 以及 substitution-aware expressions。 \\
UVec layer & \texttt{uvec.hpp}, \texttt{commands.hpp} & \texttt{UVec<T>}, pointer wrappers, handle manager & 提供 cross-device vector storage、CPU/CUDA synchronization 和 script-visible handles。 \\
Geometry CPU backend & \texttt{geometry.hpp} & \texttt{Point}, \texttt{Segment}, \texttt{BoundingBox}, batch algorithms & 在 CPU data 上实现 scalar 和 batch two-dimensional geometry algorithms。 \\
CUDA backend & \texttt{geometry\_cuda.*} & kernels and launch wrappers & 为 independent-work-item geometry algorithms 提供 GPU acceleration。 \\
Tests & \texttt{test/*.cpp} & semantic, equivalence, random, edge, stress tests & 验证 language semantics、UVec consistency 和 CPU/GPU correctness。 \\
Benchmarks and figures & \texttt{docs/data}, \texttt{docs/scripts}, \texttt{docs/figures} & CSV files and Python generators & 生成 reproducible performance data 和 thesis-ready figures。 \\
\bottomrule
\end{tabularx}
\end{table}

\begin{table}[htbp]
\centering
\caption{实现规模统计（Quantitative implementation statistics）}
\label{tab:implementation-statistics}
\small
\begin{tabularx}{\textwidth}{lrX}
\toprule
Metric & Value & Meaning \\
\midrule
C++/CUDA source files & 56 & \texttt{include}、\texttt{src} 和 \texttt{test} 下的 header、source、CUDA、benchmark 和 test files。 \\
C++/CUDA source lines & 19151 & \texttt{.hpp}、\texttt{.cpp} 和 \texttt{.cu} 文件总行数。 \\
Header/interface lines & 12881 & 核心 interpreter、UVec、parser、expression、command 和 geometry interfaces。 \\
CUDA/source implementation lines & 642 & CUDA backend 和非 header source files。 \\
Test source lines & 5628 & Regression、equivalence、edge-case、random、stress 和 benchmark-related test code。 \\
CTest entries & 31 & Build-system registered tests，用于 regression 和 validation。 \\
CUDA kernels & 17 & CUDA geometry backend 以及相关 test path 中的 device kernels。 \\
Benchmark executables & 6 & 用于 runtime、geometry 和 multi-GPU experiments 的 non-interactive benchmark programs。 \\
Benchmark/figure scripts & 5 & \texttt{docs/scripts} 下的数据和图表生成脚本。 \\
Generated SVG figures & 18 & \texttt{docs/figures} 下的 reproducible vector figures。 \\
Thesis PNG figures & 18 & 复制到 \texttt{contents/figures} 中用于论文编译的 rendered figures。 \\
\bottomrule
\end{tabularx}
\end{table}

\section{Interpreter Runtime}
Runtime 由 values、words、scripts、scopes 和 commands 组成。Value 表示脚本层的值；Word 表示一个 command argument 在求值前后的结构；Script 是 command list；Interpreter 维护 command tables 和 variable scopes。FTCL 的执行路径保持 Tcl-style 的简单性，但在内部为后续扩展保留了结构化表示。

一个普通脚本的执行过程如下：

\begin{enumerate}
  \item 将 source text 解析为 script structure。
  \item 按顺序执行一个个 command。
  \item 对每个 word 执行 variable substitution、command substitution、array lookup 和 string concatenation。
  \item 将第一个求值后的 word 作为 command name。
  \item 调用对应的 native C++ command function，并返回 value 或 exception。
\end{enumerate}

这种结构有一个重要优势：新能力不需要发明新语法。\texttt{thread}、\texttt{uvec} 和 \texttt{geom} 都是 command，而不是语言核心的特殊分支。因此 FTCL 可以保持小而清晰，同时承载并发、跨设备内存和几何算法。

\section{词法分析流程（Lexical Analysis Flow）}
Lexical layer 将原始 script string 转换为 token stream。FTCL 的 token-stream 方向会保留 token text 和 token boundaries，使后续阶段能明确识别 command endings、comments、whitespace、braces、brackets、variable references 和 ordinary text。算法 \ref{alg:lexer} 描述了 lexer 的基本流程。

\begin{algorithm}[H]
\caption{FTCL Lexical Analysis}
\label{alg:lexer}
\KwInput{source string $source$}
\KwOutput{token stream $tokens$}
$tokens \leftarrow [\ ]$\;
$cursor \leftarrow 0$\;
\While{$cursor < |source|$}{
  $ch \leftarrow source[cursor]$\;
  \uIf{$ch \in \{\mathrm{space},\mathrm{tab}\}$}{
    $start \leftarrow cursor$\;
    consume horizontal whitespace\;
    append \textsc{Token}(\textsc{Whitespace}, $source[start:cursor]$) to $tokens$\;
  }
  \uElseIf{$ch$ is a newline or $ch = \texttt{;}$}{
    append \textsc{Token}(\textsc{CommandEnd}, $ch$) to $tokens$\;
    $cursor \leftarrow cursor + 1$\;
  }
  \uElseIf{$ch = \texttt{\#}$ and $cursor$ is at command start}{
    $start \leftarrow cursor$\;
    consume until newline or end of file\;
    append \textsc{Token}(\textsc{Comment}, $source[start:cursor]$) to $tokens$\;
  }
  \uElseIf{$ch$ is a Tcl special delimiter}{
    append \textsc{Token}(\textsc{Special}, $ch$) to $tokens$\;
    $cursor \leftarrow cursor + 1$\;
  }
  \Else{
    $start \leftarrow cursor$\;
    consume ordinary characters until a token boundary\;
    append \textsc{Token}(\textsc{Text}, $source[start:cursor]$) to $tokens$\;
  }
}
append \textsc{Token}(\textsc{EndOfFile}, empty string) to $tokens$\;
\KwReturn{$tokens$}\;
\end{algorithm}

Lexer 不执行 substitution。它只记录输入中可观察的结构。这样做非常重要，因为 Tcl-style substitution 本质上是 semantic action：\texttt{\$} 可能需要 variable lookup，\texttt{[...]} 可能需要 recursive command evaluation。把 tokenization 和 evaluation 分开，可以让 parser 更稳定，也便于错误定位和回归测试。

\section{语法分析流程（Syntax Analysis Flow）}
Parser 消费 token stream，并把 tokens 分组成 commands 和 words。Command 在 newline、semicolon 或 end-of-file 处结束，但如果 parser 位于 braces、quotes 或 brackets 内部，则这些符号不再表示 command boundary。Parser 不执行 command，它只构造后续可执行的 script representation。

\begin{algorithm}[H]
\caption{Token-stream Command Parsing}
\label{alg:parse-script}
\KwInput{token stream $tokens$}
\KwOutput{script object $script$}
\Proc{\textsc{ParseScript}($tokens$)}{
  $cursor \leftarrow \textsc{TokenCursor}(tokens)$\;
  $script \leftarrow [\ ]$\;
  \While{not $cursor.atEnd()$}{
    $cursor.skipWhitespaceAndComments()$\;
    \If{$cursor.atCommandEnd()$}{
      $cursor.next()$\;
      \textbf{continue}\;
    }
    $command \leftarrow \textsc{ParseCommand}(cursor)$\;
    append $command$ to $script.commands$\;
    \uIf{$cursor.atCommandEnd()$}{
      $cursor.next()$\;
    }
    \uElseIf{not $cursor.atEnd()$}{
      raise ``expected command boundary''\;
    }
  }
  \KwReturn{$script$}\;
}
\Proc{\textsc{ParseCommand}($cursor$)}{
  $words \leftarrow [\ ]$\;
  \While{not $cursor.atEnd()$ and not $cursor.atCommandEnd()$}{
    $cursor.skipInlineWhitespace()$\;
    \If{$cursor.atCommandEnd()$ or $cursor.atEnd()$}{
      \textbf{break}\;
    }
    $word \leftarrow \textsc{ParseWord}(cursor)$\;
    append $word$ to $words$\;
  }
  \KwReturn{\textsc{Command}($words$)}\;
}
\end{algorithm}

Word parser 比 command parser 更细，因为一个 word 内部可能包含 literal text、variable references、command substitutions、array references、braced literals、quoted strings 和 expansion syntax。关键设计是：parser 将这些组成部分保存为结构化 word nodes，而 interpreter 负责决定它们在运行时的 value。

\begin{algorithm}[H]
\caption{Word Parsing}
\label{alg:parse-word}
\KwInput{token cursor $cursor$}
\KwOutput{word node}
\Proc{\textsc{ParseWord}($cursor$)}{
  \uIf{$cursor$ matches an open brace}{
    \KwReturn{\textsc{Word.Value}(\textsc{ReadBalancedBracedText}($cursor$))}\;
  }
  \uElseIf{$cursor$ matches a quote}{
    $parts \leftarrow [\ ]$\;
    \While{$cursor$ does not match the closing quote}{
      append \textsc{ParseWordPart}($cursor$, true) to $parts$\;
    }
    \KwReturn{\textsc{Word.Tokens}($parts$)}\;
  }
  \uElseIf{$cursor$ matches the expansion prefix \texttt{\{*\}}}{
    $child \leftarrow \textsc{ParseWord}(cursor)$\;
    \KwReturn{\textsc{Word.Expand}($child$)}\;
  }
  $parts \leftarrow [\ ]$\;
  \While{not $cursor.atWordBoundary()$}{
    append \textsc{ParseWordPart}($cursor$, false) to $parts$\;
  }
  \uIf{$parts$ contains one literal part}{
    \KwReturn{\textsc{Word.Value}($parts[0]$)}\;
  }
  \Else{
    \KwReturn{\textsc{Word.Tokens}($parts$)}\;
  }
}
\Proc{\textsc{ParseWordPart}($cursor$, $quoted$)}{
  \uIf{$cursor.peek() = \texttt{\$}$}{
    \KwReturn{\textsc{ParseVariableReference}($cursor$)}\;
  }
  \uElseIf{$cursor.peek() = \texttt{[}$}{
    \KwReturn{\textsc{Word.Script}(\textsc{ReadBalancedBracketScript}($cursor$))}\;
  }
  \uElseIf{$cursor.peek()$ is a backslash}{
    \KwReturn{\textsc{Word.String}(\textsc{ParseBackslashSequence}($cursor$))}\;
  }
  \Else{
    \KwReturn{\textsc{Word.String}(\textsc{ReadLiteralText}($cursor$, $quoted$))}\;
  }
}
\end{algorithm}

\section{语义执行流程（Semantic Execution Flow）}
Semantic execution 发生在 parsing 之后。Interpreter 逐条执行 command，对每个 word 求值，并把第一个 value 作为 command name，剩余 values 作为 arguments。Command function 才真正改变变量、创建 UVec、启动 CUDA kernel 或返回结果。

\begin{algorithm}[H]
\caption{Command-level Semantic Execution}
\label{alg:semantic-exec}
\KwInput{interpreter $interp$, parsed script $script$}
\KwOutput{last command result}
\Proc{\textsc{EvalScript}($interp$, $script$)}{
  $result \leftarrow$ empty string\;
  \ForEach{$command \in script.commands$}{
    \If{$command.words$ is empty}{
      \textbf{continue}\;
    }
    $values \leftarrow [\ ]$\;
    \ForEach{$word \in command.words$}{
      $value \leftarrow \textsc{EvalWord}(interp, word)$\;
      \uIf{$word$ is an expansion word}{
        extend $values$ by \textsc{ParseList}($value$)\;
      }
      \Else{
        append $value$ to $values$\;
      }
    }
    $name \leftarrow values[0].asString()$\;
    $proc \leftarrow interp.lookupCommand(name)$\;
    \If{$proc$ is missing}{
      raise ``invalid command name''\;
    }
    $result \leftarrow proc(interp, values)$\;
    \If{$result.code$ is \textsc{Return}, \textsc{Break}, \textsc{Continue}, or \textsc{Error}}{
      propagate $result$ according to the current context\;
    }
  }
  \KwReturn{$result$}\;
}
\end{algorithm}

\begin{algorithm}[H]
\caption{Word Evaluation}
\label{alg:word-eval}
\KwInput{interpreter $interp$, word node $word$}
\KwOutput{runtime value}
\Proc{\textsc{EvalWord}($interp$, $word$)}{
  \uIf{$word.kind = \textsc{Value}$}{
    \KwReturn{$word.literal$}\;
  }
  \uElseIf{$word.kind = \textsc{String}$}{
    \KwReturn{$word.text$}\;
  }
  \uElseIf{$word.kind = \textsc{VarRef}$}{
    \KwReturn{$interp.lookupVariable(word.name)$}\;
  }
  \uElseIf{$word.kind = \textsc{ArrayRef}$}{
    $index \leftarrow \textsc{EvalWord}(interp, word.index)$\;
    \KwReturn{$interp.lookupArrayElement(word.name, index)$}\;
  }
  \uElseIf{$word.kind = \textsc{Script}$}{
    \KwReturn{\textsc{EvalScript}($interp$, \textsc{ParseScript}(\textsc{Tokenize}($word.script$)))}\;
  }
  \Else{
    $output \leftarrow$ empty string\;
    \ForEach{$part \in word.parts$}{
      $output \leftarrow output + \textsc{EvalWord}(interp, part).asString()$\;
    }
    \KwReturn{$output$}\;
  }
}
\end{algorithm}

\section{UVec 跨设备设计}
UVec 是一个 vector-like container，但它不只拥有 CPU storage。普通 \texttt{std::vector<T>} 只维护 CPU memory；UVec 同时维护 CPU buffer 和 optional CUDA buffers，并记录哪一个 device 上的数据是 valid and up-to-date。

当 CPU code 请求最新数据，而最新副本在 CUDA 上时，UVec 执行 CUDA-to-CPU copy；当 CUDA kernel 请求最新数据，而最新副本在 CPU 上时，UVec 执行 CPU-to-CUDA copy；当某个 device 请求 mutable pointer 时，其他副本会被标记为 invalid，因为该 device 可能写入新值。图 \ref{fig:uvec-sync} 展示同步机制，图 \ref{fig:uvec-state} 展示状态机视角。

\begin{figure}[H]
  \centering
  \includegraphics[width=0.88\textwidth]{uvec-sync.png}
  \caption{UVec 跨设备同步机制（Cross-device synchronization in UVec）}
  \label{fig:uvec-sync}
\end{figure}

\begin{figure}[H]
  \centering
  \includegraphics[width=0.75\textwidth]{uvec-state.png}
  \caption{UVec 有效性状态机（UVec validity-state machine）}
  \label{fig:uvec-state}
\end{figure}

最重要的两个 pointer interfaces 是 \texttt{as\_uptr(device)} 和 \texttt{as\_mut\_uptr(device)}。前者返回 read-only pointer，并保证目标 device 上的数据已经 valid；后者返回 mutable pointer，并将目标 device 标记为 newest writer。这两个接口让 CPU algorithm 和 CUDA kernel 可以共享同一个数据抽象，而不需要脚本用户手动判断数据在哪里。

\section{形式化 UVec 一致性模型}
UVec 可以被描述为定义在 device set 上的小型 consistency model。令
\[
  D = \{\mathrm{CPU}, \mathrm{CUDA}_0, \mathrm{CUDA}_1, \ldots\}
\]
表示可能的存储设备集合。长度为 \(n\) 的 UVec instance 表示为
\[
  U = (n, B, V, L),
\]
其中 \(B_d\) 是 device \(d \in D\) 上的 physical buffer，\(V_d \in \{0,1\}\) 表示 \(B_d\) 是否包含 valid copy，\(L \subseteq D\) 表示可能包含 latest value 的 device set。当前实现通常表现为 single latest writer，但集合表示允许 CPU 和 CUDA 在一次 read synchronization 后同时 valid。

两个核心操作满足以下规则：

\begin{description}
  \item[Read rule:] \texttt{as\_uptr(d)} 只有在保证 \(V_d = 1\) 后才能返回 device \(d\) 上的 pointer。如果 \(V_d = 0\)，则从某个满足 \(V_s = 1\) 的 source device \(s\) 复制数据，并设置 \(V_d \leftarrow 1\)。Read operation 不会使其他 valid copies 失效。
  \item[Write rule:] \texttt{as\_mut\_uptr(d)} 返回 device \(d\) 上的 mutable pointer。返回前必须保证 \(B_d\) 存在且包含 latest value；随后设置 \(V_d \leftarrow 1\)，并对所有 \(s \ne d\) 设置 \(V_s \leftarrow 0\)，因为通过 mutable pointer 的后续写入可能改变 logical vector。
\end{description}

由此得到两个有用的不变式：
\[
  \exists d \in D,\; V_d = 1
\]
在每次成功初始化后的 UVec operation 之后成立；并且 UVec 返回的每个 read pointer 在返回时都指向 logical vector 的 valid copy。这个模型刻意保持简单：它足以支撑同步式 CPU/CUDA command execution，同时也为未来的 multiple CUDA devices、asynchronous streams 和 copy-on-write policies 留出空间。

\section{Geometry Command Flow}
\texttt{geom} command 的执行流程如图 \ref{fig:geom-flow} 所示。以 \texttt{geom nearest\_point} 为例，interpreter 首先求值 script arguments；command 层解析 UVec handles 和 device options；如果目标是 CPU，就把输入同步到 CPU 并调用 C++ algorithm；如果目标是 CUDA，就把输入同步到 CUDA，创建 output UVec，启动 CUDA kernel，并返回 output handle。

\begin{figure}[H]
  \centering
  \includegraphics[width=0.90\textwidth]{geom-flow.png}
  \caption{geom 命令执行流程（Execution flow of a geometry command）}
  \label{fig:geom-flow}
\end{figure}

Handle 机制非常关键。大规模点集不应该在每条命令之间反复转换为 script strings。更合理的方式是：script 创建一次 UVec，把 handle 传给多个 geometry commands，只在最终需要展示或统计时再 read back。这个思路也是后面 performance evaluation 中 speedup 能出现的前提。

\section{并发设计}
FTCL 提供 thread tasks 和 thread channels。Thread task 可以异步运行一段 script，并在之后被 await；thread channel 是由 mutex 和 condition variable 保护的 message queue，支持 blocking receive 和 non-blocking try-receive。

这个 concurrency model 对实验很有用。它可以驱动 game-like demo，模拟 producer-consumer workload，也可以通过多个 worker scripts 分发 geometry requests。当前 CUDA concurrency experiment 也暴露出一个限制：多个 script workers 可能竞争同一个 GPU，所以未来需要 streams 或 dedicated GPU scheduler。

后续第 5 章会分别从 single-GPU contention 和 multi-GPU weak-scaling 两个角度评价这个设计。这些实验也解释了为什么实现中需要把 UVec handle manager 从 global lock 细化为 handle-level locking。

\section{小结}
本章给出了 FTCL 的系统设计：语言层负责表达，command 层负责连接，UVec 层负责跨设备数据一致性，CPU/CUDA backend 负责具体计算。这个架构的关键不在于某一个单独模块，而在于这些模块之间的边界足够清楚，因而能够支持后续测试、benchmark 和应用场景扩展。

FTCL separates control, data residency, and execution backend. UVec is the bridge that keeps script-level geometry commands independent from low-level CPU/CUDA memory management.

\chapter{Implementation}

\section{Interpreter Core}
The FTCL implementation is mainly header-based and located under the \texttt{include} directory. The interpreter is built as a CMake interface library, while tests and benchmark executables are placed under the \texttt{test} directory. Optional CUDA support is enabled by the \texttt{FTCL\_ENABLE\_CUDA} CMake option.

Value representation follows the idea that script values are strings but may also cache structured internal data. For example, a value can be constructed from an integer, floating-point number, Boolean, list, dictionary, or variable name. When the string representation is requested, it is generated lazily. This reduces repeated formatting while preserving a string-oriented script interface.

\section{Expression Parser}
The expression parser implements a recursive-descent precedence hierarchy. The top-level function parses conditional expressions. It then descends through logical OR, logical AND, bitwise operators, equality, relational operators, shifts, additive operators, multiplicative operators, unary operators, and primary expressions.

For example, \texttt{a || b \&\& c} is parsed by the logical OR layer. The left operand is parsed as a logical AND expression. After the OR operator is matched, the right operand is again parsed as a logical AND expression, so \texttt{b \&\& c} is grouped before the OR operation is applied.

The parser also supports short-circuit evaluation. If the left operand of \texttt{||} is true, the right operand is parsed but not evaluated. This prevents side effects in skipped branches while preserving syntax validation and parser position.

\begin{algorithm}[H]
\caption{Expression Evaluation and Conditional Operator}
\label{alg:expr-eval}
\KwInput{interpreter $interp$, expression string $input$}
\KwOutput{expression value}
\Proc{\textsc{EvalExpr}($interp$, $input$)}{
  $parser \leftarrow \textsc{ExprParser}(interp, input)$\;
  $value \leftarrow parser.parseConditional()$\;
  $parser.skipWhitespace()$\;
  \If{not $parser.atEnd()$}{
    raise ``unexpected token in expression''\;
  }
  \KwReturn{$value$}\;
}
\Proc{\textsc{ParseConditional}($parser$)}{
  $condition \leftarrow \textsc{ParseLogicalOr}(parser)$\;
  \If{not $parser.match(\texttt{?})$}{
    \KwReturn{$condition$}\;
  }
  $takeTrue \leftarrow condition.isTrue()$\;
  $parser.setEvaluationEnabled(takeTrue)$\;
  $trueValue \leftarrow \textsc{ParseLogicalOr}(parser)$\;
  $parser.expect(\texttt{:})$\;
  $parser.setEvaluationEnabled(\neg takeTrue)$\;
  $falseValue \leftarrow \textsc{ParseConditional}(parser)$\;
  $parser.restoreEvaluationMode()$\;
  \uIf{$takeTrue$}{
    \KwReturn{$trueValue$}\;
  }
  \Else{
    \KwReturn{$falseValue$}\;
  }
}
\end{algorithm}

\begin{algorithm}[H]
\caption{Logical OR Short-circuit Evaluation}
\label{alg:expr-or}
\KwInput{expression parser $parser$}
\KwOutput{logical OR result}
\Proc{\textsc{ParseLogicalOr}($parser$)}{
  $lhs \leftarrow \textsc{ParseLogicalAnd}(parser)$\;
  \While{$parser.match(\texttt{||})$}{
    \uIf{$lhs.isTrue()$}{
      $parser.setEvaluationEnabled(false)$\;
      \textsc{ParseLogicalAnd}($parser$)\tcp*[r]{parse only; no side effects}
      $parser.restoreEvaluationMode()$\;
      $lhs \leftarrow 1$\;
    }
    \Else{
      $rhs \leftarrow \textsc{ParseLogicalAnd}(parser)$\;
      $lhs \leftarrow bool(lhs) \lor bool(rhs)$\;
    }
  }
  \KwReturn{$lhs$}\;
}
\end{algorithm}

\section{UVec Implementation}
UVec requires its element type to be trivially copyable. This is necessary because CPU/CUDA transfers are implemented as byte copies. The main supporting types are \texttt{Device}, \texttt{RawUPtr<T>}, \texttt{RawUMutPtr<T>}, \texttt{UniversalCopy}, and \texttt{Zeroable}.

The script-level \texttt{uvec} command manages UVec handles. It can create vectors from lists, create filled or zeroed vectors, get or set elements, return vector length, synchronize to a device, and convert a vector back to a script list. The command manager stores handle-to-UVec mappings under a mutex so that scripts can pass handles safely between commands.

\begin{algorithm}[H]
\caption{UVec Read Synchronization}
\label{alg:uvec-read}
\KwInput{UVec object $uvec$, target device $device$}
\KwOutput{read-only unified pointer}
\Proc{\textsc{AsUPtr}($uvec$, $device$)}{
  \uIf{$device = \mathrm{CPU}$}{
    \If{CPU buffer is not valid}{
      $source \leftarrow$ any valid CUDA buffer\;
      copy $source$ CUDA buffer to CPU buffer\;
      mark CPU buffer valid\;
    }
    \KwReturn{\textsc{RawUPtr}(CPU pointer, length, CPU)}\;
  }
  \Else{
    ensure CUDA buffer for $device$ exists\;
    \If{target CUDA buffer is not valid}{
      \uIf{CPU buffer is valid}{
        copy CPU buffer to target CUDA buffer\;
      }
      \Else{
        $source \leftarrow$ another valid CUDA buffer\;
        copy $source$ CUDA buffer to target CUDA buffer\;
      }
      mark target CUDA buffer valid\;
    }
    \KwReturn{\textsc{RawUPtr}(CUDA pointer, length, $device$)}\;
  }
}
\end{algorithm}

\begin{algorithm}[H]
\caption{UVec Mutable Pointer Access}
\label{alg:uvec-write}
\KwInput{UVec object $uvec$, target device $device$}
\KwOutput{mutable unified pointer}
\Proc{\textsc{AsMutUPtr}($uvec$, $device$)}{
  ensure storage exists on $device$\;
  \If{target device is not valid}{
    synchronize latest data to target device\;
  }
  \ForEach{$other \in$ devices and $other \ne device$}{
    mark $other$ invalid\;
  }
  mark $device$ valid and newest\;
  \KwReturn{\textsc{RawUMutPtr}(target pointer, length, $device$)}\;
}
\end{algorithm}

\section{CPU Geometry Algorithms}
The CPU geometry library is implemented in \texttt{include/geometry.hpp}. The basic types are Point, Segment, and BoundingBox. The library includes distance, squared distance, dot product, cross product, orientation, point-on-segment checks, segment intersection, line intersection, point-line distance, point-segment distance, polygon area, polygon perimeter, point-in-polygon classification, convex hull, bounding box, and closest-pair distance.

Batch algorithms use flattened arrays. A point set is stored as \(x_0,y_0,x_1,y_1,\ldots\). A segment set is stored as \(x_1,y_1,x_2,y_2\) per segment. An AABB set is stored as \(min_x,min_y,max_x,max_y\) per box. This format is easy for scripts, UVec, and CUDA kernels to share.

\begin{table}[H]
\centering
\caption{Representative geometry algorithms and asymptotic costs}
\label{tab:algorithm-complexity}
\footnotesize
\setlength{\tabcolsep}{3.5pt}
\begin{tabularx}{\textwidth}{@{}>{\ttfamily\raggedright\arraybackslash}p{0.31\textwidth} >{\raggedright\arraybackslash}X >{\centering\arraybackslash}p{0.12\textwidth} >{\raggedright\arraybackslash}X@{}}
\toprule
Command or function & Input scale & CPU complexity & GPU suitability \\
\midrule
distance & one point pair & \(O(1)\) & Low; scalar command \\
segment\_intersect & one segment pair & \(O(1)\) & Low; scalar command \\
polygon\_area & \(P\) polygon vertices & \(O(P)\) & Medium for batched polygons \\
convex\_hull & \(N\) points & \(O(N\log N)\) & Low in this project \\
bbox\_reduce & \(N\) points & \(O(N)\) & High with parallel reduction \\
centroid & \(N\) points & \(O(N)\) & High with parallel reduction \\
batch\_distance\_matrix & \(N\) by \(Q\) points & \(O(NQ)\) & Very high; independent outputs \\
nearest\_point & \(N\) data, \(Q\) queries & \(O(NQ)\) & High; independent queries \\
range\_count\_circle & \(N\) points, \(Q\) circles & \(O(NQ)\) & High; independent queries \\
batch\_segment\_intersect & \(K\) segment pairs & \(O(K)\) & High; independent pairs \\
collision\_aabb & \(K\) box pairs & \(O(K)\) & High; independent pairs \\
transform\_points & \(N\) points & \(O(N)\) & Very high; independent points \\
\bottomrule
\end{tabularx}
\end{table}

The table explains why not every geometry command is accelerated in the same way. Scalar commands are useful for script expressiveness but are too small to amortize CUDA launch overhead. Batch commands with \(O(NQ)\) or \(O(N)\) independent work items are more suitable for GPU execution. Algorithms such as convex hull require sorting and sequential stack-like structure, so they are kept as CPU-first algorithms in this thesis.

\section{CUDA Geometry Backend}
The CUDA backend is implemented in \texttt{src/geometry\_cuda.cu}.
It is declared in \texttt{include/geometry\_cuda.hpp}.
It supports data-parallel algorithms such as distance matrix, nearest point, circle range count, point transformation, batch orientation, AABB collision, polygon batch area, and distance to polyline.

For a distance matrix with \(N\) source points and \(Q\) query points, the output has \(NQ\) entries. A CUDA thread computes one output entry by loading one source point and one query point, computing the Euclidean distance, and writing the result. This mapping is simple and effective because each output element is independent.

\section{Script Interface}
The following example shows how a script can create two point sets, run a CUDA nearest-point query, and read back the result.

\begin{lstlisting}[language=bash,caption={FTCL geometry scripting example}]
set dataset [geom uvec_points {{0 0} {1 2} {3 4} {5 8}} cuda:0]
set query   [geom uvec_points {{2 2} {4 6}} cuda:0]
set result  [geom nearest_point $dataset $query cuda:0]
puts [uvec to_list $result]
\end{lstlisting}

The user does not manually allocate device memory. The script describes the task, UVec manages synchronization, and the geom command selects the backend.

\begin{algorithm}[H]
\caption{Geometry Command Dispatch for Nearest-point Query}
\label{alg:geom-dispatch}
\KwInput{interpreter $interp$, command arguments $argv$}
\KwOutput{UVec handle containing nearest-point results}
\Proc{\textsc{GeomNearestPoint}($interp$, $argv$)}{
  $datasetId \leftarrow \textsc{ParseHandle}(argv[2])$\;
  $queryId \leftarrow \textsc{ParseHandle}(argv[3])$\;
  $device \leftarrow \textsc{ParseOptionalDevice}(argv[4], \mathrm{CPU})$\;
  $dataset \leftarrow uvecManager.lookup(datasetId)$\;
  $queries \leftarrow uvecManager.lookup(queryId)$\;
  \uIf{$device$ is CUDA}{
    $datasetPtr \leftarrow dataset.asUPtr(device)$\;
    $queryPtr \leftarrow queries.asUPtr(device)$\;
    $output \leftarrow \textsc{UVec.Uninitialized}(queryCount \times 2, device)$\;
    $outputPtr \leftarrow output.asMutUPtr(device)$\;
    launch \textsc{NearestPointCuda}($outputPtr$, $datasetPtr$, $queryPtr$)\;
    \KwReturn{$uvecManager.create(output)$}\;
  }
  \Else{
    $datasetCpu \leftarrow dataset.cpuVector()$\;
    $queriesCpu \leftarrow queries.cpuVector()$\;
    $values \leftarrow \textsc{NearestPointCpu}(datasetCpu, queriesCpu)$\;
    $output \leftarrow \textsc{UVec.FromCpu}(values)$\;
    \KwReturn{$uvecManager.create(output)$}\;
  }
}
\end{algorithm}

\chapter{实验评估}

\section{Evaluation Goals}
本章实验回答三个问题。第一，FTCL 的 language runtime 是否保持了预期语义？第二，CPU 和 CUDA geometry backends 是否能产生等价结果？第三，从 FTCL script 用户的角度看，CUDA 在什么规模下能够带来 end-to-end speedup？

项目的实验数据、CSV 文件和图表生成脚本都位于 \texttt{docs} 目录下。测量内容包括 semantic pass rate、parser backend timing、thread-channel latency、frame-time distribution、CPU/CUDA geometry speedup、throughput、numerical error、break-even point、ablation data 和 concurrent geometry performance。这些指标共同覆盖 correctness、performance 和 scalability 三个层面。

\section{Experimental Methodology}
除非特别说明，本文所有 performance results 都被视为 end-to-end system measurements。一次 geometry command measurement 包含 FTCL command dispatch、argument parsing、UVec handle lookup、device synchronization、output UVec allocation、backend execution，以及 CUDA kernel launch 后的 synchronization。它不是纯 CUDA kernel time。本文有意采用这种测量方式，因为论文评价的是脚本用户实际看到的系统性能。

Benchmark pipeline 遵循以下原则：

\begin{enumerate}
  \item \textbf{使用 Release builds。} Debug builds 不进入性能讨论，因为 interpreter dispatch 和 STL containers 对 optimization level 非常敏感。
  \item \textbf{区分 warmup 和 measurement。} Parser backend timing 和 geometry timing scripts 在可支持的情况下会先执行 warmup runs，再记录 samples。
  \item \textbf{优先使用 median 或 percentile。} Latency distributions 使用 P50、P95 和 P99 描述；geometry timing 在有多个样本时使用 representative median command latency。
  \item \textbf{单独看待 readback cost。} 部分实验将结果保留在 UVec handles 中，ablation experiments 则显式测量将 large UVec outputs 转回 script lists 的成本。
  \item \textbf{先比较正确性，再解释速度。} 只有 CPU/GPU equivalence tests 在数值容忍范围内通过时，CUDA speedups 才有意义。
\end{enumerate}

主要派生指标定义如下：
\[
  \mathrm{speedup} = \frac{T_{\mathrm{CPU}}}{T_{\mathrm{CUDA}}},
  \quad
  \mathrm{throughput} = \frac{\mathrm{work\ items}}{T},
  \quad
  \mathrm{work\ items} = N \times Q.
\]
对于 distance matrix、nearest point 和 range-count workloads，\(N\) 表示 data points 数量，\(Q\) 表示 query points 或 query regions 数量。

\subsection{CPU Baseline Definition}
本文的 CPU baseline 是 FTCL CPU geometry backend，而不是手工优化的 parallel CPU library。这个区别很重要：实验回答的问题是脚本用户会问的问题，即同一个 FTCL command 在 CPU device 和 CUDA device 上执行时，end-to-end improvement 是多少。

因此，CPU row 和 CUDA row 尽量共享相同的 script-level command path。主性能实验中，输入会在计时前预先构造为 UVec handles；计时区间包括 command dispatch、handle lookup、output UVec allocation、backend execution 和必要 synchronization。CPU row 使用 C++ scalar loops 处理 CPU-resident UVec data；CUDA row 使用 CUDA device buffers 和 kernels，并且当前 CUDA helper 在 kernel launch 后调用 \texttt{cudaDeviceSynchronize}，保证 command 返回时 backend work 已完成。

这个 baseline 对评价 FTCL 作为 heterogeneous scripting system 是公平的，但不能被解释为对抗高度优化的 all-core CPU geometry library。实验服务器拥有 128 个 logical CPU threads，而当前 CPU geometry backend 故意保持简单。因此，报告的 speedup 衡量的是在相同 command interface 下，将 FTCL backend 从 scalar CPU implementation 移动到 CUDA implementation 后获得的收益。未来加入 all-core CPU backend 会是有价值的补充 baseline。

\begin{table}[H]
\centering
\caption{实验环境信息（Experimental environment for reproducibility）}
\label{tab:environment}
\begin{tabularx}{\textwidth}{lX}
\toprule
Item & Recorded value \\
\midrule
Operating system & WSL Ubuntu environment used for build, tests, CUDA execution, and figure generation \\
CPU model & Intel Xeon Platinum 8358 CPU at 2.60 GHz \\
CPU topology & 2 sockets, 32 cores per socket, 2 hardware threads per core, 128 logical CPUs \\
CPU frequency & 800 MHz minimum, 3.40 GHz maximum as reported by \texttt{lscpu} \\
NUMA layout & 2 NUMA nodes; node 0 contains CPUs 0--31 and 64--95, node 1 contains CPUs 32--63 and 96--127 \\
Main memory & 1.0 TiB total memory; about 623 GiB available when the environment was recorded \\
GPU model & 8 NVIDIA A800-SXM4-80GB GPUs \\
GPU memory & 81920 MiB per GPU, approximately 640 GiB aggregate device memory \\
GPU driver & NVIDIA driver 530.30.02 \\
CUDA compute capability & 8.0 for each A800 GPU \\
Build system & CMake with Release configuration \\
C++ standard & C++20, selected by \texttt{FTCL\_CXX\_STANDARD=20} \\
Compiler & GCC or Clang available inside WSL \\
CUDA option & \texttt{FTCL\_ENABLE\_CUDA=ON} for CUDA geometry experiments \\
CPU build path & \texttt{build} or \texttt{/tmp/ftcl-build-geometry-cpu} \\
CUDA build path & \texttt{/tmp/ftcl-build-geometry-cuda} \\
Benchmark data & CSV files under \texttt{docs/data/benchmarks} and \texttt{docs/data/geometry} \\
Generated figures & SVG/PNG figures under \texttt{docs/figures} and \texttt{contents/figures} \\
Important note & Free GPU memory is time-dependent because other users or jobs may share the same server \\
\bottomrule
\end{tabularx}
\end{table}

表 \ref{tab:environment} 是环境记录，而不是硬件无关结论。本文所有 performance numbers 都应理解为在该 multi-GPU server 上的实测结果。服务器的 CPU core count 也说明 CPU baseline 并不来自弱硬件；真正的实验边界在于当前 CPU backend 采用 scalar implementation，而不是 fully parallel CPU implementation。

\section{Semantic Regression Tests}
Semantic regression tests 覆盖 core commands、collections、control flow、expressions 和 threading。图 \ref{fig:semantic-pass} 显示当前测量的 test suites 达到 100 percent pass rate。

\begin{figure}[H]
  \centering
  \includegraphics[width=0.78\textwidth]{semantic-pass-rate.png}
  \caption{测试套件语义通过率（Semantic pass rate by test suite）}
  \label{fig:semantic-pass}
\end{figure}

这个结果并不表示 FTCL 已经完整实现 Tcl。它表示本文选定的 Tcl subset 已经形成稳定 regression baseline。对于 parser 和 expression 这类容易引入细微语义变化的模块，这一点尤其重要。

\section{Parser Backend Timing}
图 \ref{fig:parser-backend} 比较 legacy parser 和 token-stream parser 在 Tcl subset benchmark 上的耗时。Token-stream path 并不一定更快，因为它增加了 lexing pass，并保存 token metadata。它的价值主要是架构层面的：更好的 diagnostics、更清晰的 command boundaries，以及通向 AST 和 bytecode execution 的路径。

\begin{figure}[H]
  \centering
  \includegraphics[width=0.80\textwidth]{parser-backend-timing.png}
  \caption{Parser backend 耗时对比（Timing comparison between parser backends）}
  \label{fig:parser-backend}
\end{figure}

直接结论：token-stream parser 的短期意义不是速度，而是可维护性、错误定位和后续编译路径。

\section{Concurrency and Runtime Smoothness}
Thread-channel latency 和 frame-time distribution 如图 \ref{fig:channel-latency} 与图 \ref{fig:frame-time} 所示。Channel experiment 通过 ping-pong round trip 估计 one-way latency；frame-time experiment 测量一个 non-interactive game-like workload。

\begin{figure}[H]
  \centering
  \includegraphics[width=0.78\textwidth]{channel-latency.png}
  \caption{线程通道单向延迟分布（Thread-channel one-way latency distribution）}
  \label{fig:channel-latency}
\end{figure}

\begin{figure}[H]
  \centering
  \includegraphics[width=0.78\textwidth]{frame-time.png}
  \caption{类游戏 workload 的帧耗时分布（Frame-time distribution）}
  \label{fig:frame-time}
\end{figure}

这些实验说明 FTCL 可以支持简单 concurrent scripts 和 interactive demos。相比平均值，P50、P95 和 P99 这类 tail-latency metrics 更能反映脚本系统的实际交互体验。

\section{CPU/GPU Equivalence}
对 CUDA backend 而言，correctness 比 speed 更重要。Equivalence tests 会让同一 geometry workload 同时在 CPU 和 CUDA 上运行，比较返回的 UVec values，并允许较小 floating-point tolerance。图 \ref{fig:error-dist} 展示 numerical error distribution。

\begin{figure}[H]
  \centering
  \includegraphics[width=0.70\textwidth]{geometry-error-distribution.png}
  \caption{CPU 与 CUDA 几何结果的数值误差分布（Numerical error distribution）}
  \label{fig:error-dist}
\end{figure}

生成数据中的 observed errors 接近零。这支持了本文的基本正确性主张：在测试覆盖到的 workloads 上，CUDA implementation 保持了 CPU geometry semantics。

\section{Geometry Performance Scaling}
性能实验聚焦三个代表性算法：distance matrix、nearest point 和 circle range count。它们的共同点是存在大量 independent work items。表 \ref{tab:geom-speedup} 根据保存的 performance CSV 总结 end-to-end speedup。

Single-GPU performance-scale experiment 使用一个 FTCL interpreter 和一个 CUDA device。每个 scale 下，benchmark 会在计时前构造 CPU 和 CUDA UVec inputs，执行可选 untimed warmup，并报告 median command time。这意味着 input point-list parsing cost 被排除在主计时之外，而 output allocation 和 backend synchronization 仍然包含在内。实验目标是测量 prebuilt UVec handles 在 realistic script pipeline 中的使用成本。

\begin{table}[H]
\centering
\caption{代表性几何算法的 end-to-end CPU/CUDA speedup}
\label{tab:geom-speedup}
\begin{tabular}{lrrr}
\toprule
Algorithm & N,Q & Work items & Speedup \\
\midrule
Distance matrix & 2048, 256 & 524288 & 6.77x \\
Nearest point & 2048, 256 & 524288 & 2.05x \\
Circle range count & 2048, 256 & 524288 & 1.41x \\
Distance matrix & 8192, 1024 & 8388608 & 51.90x \\
Nearest point & 8192, 1024 & 8388608 & 8.39x \\
Circle range count & 8192, 1024 & 8388608 & 6.50x \\
Distance matrix & 32768, 2048 & 67108864 & 224.00x \\
Nearest point & 32768, 2048 & 67108864 & 16.55x \\
Circle range count & 32768, 2048 & 67108864 & 13.13x \\
\bottomrule
\end{tabular}
\end{table}

\begin{figure}[H]
  \centering
  \includegraphics[width=0.82\textwidth]{geometry-speedup.png}
  \caption{批量几何算法的 CPU/CUDA speedup}
  \label{fig:geom-speedup}
\end{figure}

图 \ref{fig:execution-bars} 和图 \ref{fig:throughput} 展示 execution time 和 throughput。Distance matrix 获得最高 speedup，因为它的 work items 简单且彼此独立。Nearest point 和 range count 同样受益于 CUDA，但每个 query 需要更多 memory reads，因此 scaling pattern 不完全相同。

\begin{figure}[H]
  \centering
  \includegraphics[width=0.82\textwidth]{geometry-execution-bars.png}
  \caption{最大测量规模下的执行时间对比（Execution-time comparison）}
  \label{fig:execution-bars}
\end{figure}

\begin{figure}[H]
  \centering
  \includegraphics[width=0.82\textwidth]{geometry-throughput.png}
  \caption{几何算法吞吐量曲线（Throughput curve）}
  \label{fig:throughput}
\end{figure}

直接结论：work items 足够多且输出紧凑时，CUDA backend 的优势更稳定；输出规模过大时，allocation 和 memory bandwidth 会成为瓶颈。

\section{Multi-GPU Scaling}
实验服务器包含 8 张 A800 GPUs，因此本文进一步测量 FTCL 是否真的能把多个 CUDA devices 暴露给 script-level geometry workloads。Multi-GPU benchmark 使用 weak scaling：每张 GPU 接收相同数量的 query points，而总 query count 随 GPU 数量增长。每个 worker 拥有独立 FTCL interpreter，在不同 CUDA device 上创建 UVec inputs，并分别在 \texttt{cuda:0}、\texttt{cuda:1}、\ldots、\texttt{cuda:7} 上启动相同 geometry command。报告的 speedup 是相对于 1 GPU 的 throughput speedup。

Benchmark 运行时需要让所有 devices 可见，例如：

\begin{lstlisting}[language=bash]
CUDA_VISIBLE_DEVICES=0,1,2,3,4,5,6,7 \
FTCL_MULTI_GPU_MODE=paper \
FTCL_MULTI_GPU_SCALING=weak \
test/bench_geometry_multi_gpu
\end{lstlisting}

记录运行中，\(N=32768\)，每张 GPU 接收 \(Q=4096\) 个 query objects，每行报告三轮 median。因此 1、2、4、8 GPUs 对应的 total query count 分别是 \(4096\)、\(8192\)、\(16384\) 和 \(32768\)。Timer 测量从同时释放所有 worker threads 到所有 workers 获得 output UVec handles 的 wall-clock interval。Length checking 和 output destruction 在计时区间之后执行，因此计时聚焦 command dispatch、UVec handle lookup、output allocation、CUDA launch、backend synchronization 和 result-handle creation。

这个实验还暴露了一个重要 runtime bottleneck。原始 UVec command manager 使用一个 global mutex 保护整个 handle table，并在持锁状态下执行 geometry commands。这会导致彼此独立的 CUDA commands 在 script runtime layer 被串行化。实现后来改为：global mutex 只保护 handle-table lookup，每个 UVec handle 再拥有独立 entry mutex。这样，不同 GPU 上的独立 UVec handles 可以并发执行。

\begin{table}[H]
\centering
\caption{多 GPU 几何命令 weak-scaling speedup}
\label{tab:multi-gpu-speedup}
\begin{tabular}{lrrrr}
\toprule
Algorithm & 1 GPU & 2 GPUs & 4 GPUs & 8 GPUs \\
\midrule
Distance matrix & 1.00x & 0.96x & 0.97x & 1.12x \\
Nearest point & 1.00x & 2.00x & 3.99x & 3.74x \\
Circle range count & 1.00x & 1.98x & 3.97x & 7.90x \\
\bottomrule
\end{tabular}
\end{table}

\begin{figure}[H]
  \centering
  \includegraphics[width=0.82\textwidth]{geometry-multi-gpu-scaling.png}
  \caption{1、2、4、8 张 A800 GPU 的 weak-scaling speedup}
  \label{fig:multi-gpu-scaling}
\end{figure}

直接结论：circle range count 接近 linear multi-GPU throughput scaling，而 distance matrix 暴露出 output-dominant bottleneck。

结果说明 multi-GPU execution 有价值，但强烈依赖 algorithm shape。Circle range counting 几乎线性扩展到 8 GPUs，因为每个 query 只产生一个 compact result，并且 work pattern 规则。Nearest point 扩展到 4 GPUs 表现良好，但到 8 GPUs 后下降，说明 runtime overhead、allocation 和 synchronization 开始占据主导。Distance matrix 的 multi-GPU scaling 最弱，因为每个 query 会产生 large output row；output allocation 和 global memory traffic 主导了 command time。这是一个有价值的 negative result：系统已经可以使用 8 GPUs，但要让 large-output kernels 真正扩展，还需要 memory pools、persistent output buffers 和 CUDA streams。

\section{Break-even Analysis}
CUDA acceleration 存在 fixed overhead。小规模 workload 可能在 CUDA 上更慢，因为 kernel launch、synchronization、allocation 和 command dispatch 会压过真正的计算量。图 \ref{fig:break-even} 使用 work items \(N \times Q\) 作为横轴，speedup 作为纵轴。

\begin{figure}[H]
  \centering
  \includegraphics[width=0.82\textwidth]{geometry-break-even.png}
  \caption{按 work items 统计的 CUDA break-even analysis}
  \label{fig:break-even}
\end{figure}

直接结论：只有 work-item count 足以摊销 launch、allocation、synchronization 和 UVec management costs 后，CUDA 才真正值得使用。

保存的数据估计 distance matrix 的 break-even point 约为 8192 work items，nearest point 约为 47848，circle range count 约为 121209。这些值不是 universal constants。它们依赖 hardware、driver state、input size、output size 以及 measurement 是否包含 readback。

\section{Ablation and Concurrent Geometry}
Ablation experiment 将 inline literal parsing、prebuilt UVec handles 和 readback overhead 分离。实验显示，prebuilt handles 可以显著降低 command latency，尤其对 CUDA commands 更明显。同时，将 large outputs 转回 script list 可能支配总时间。因此，高性能 FTCL script 应尽量让数据在 UVec form 中跨多个 geometry commands 保持驻留，只在最后读取 summary results。

\begin{table}[H]
\centering
\caption{几何消融实验解释（Interpretation of geometry ablation experiment）}
\label{tab:ablation-interpretation}
\begin{tabularx}{\textwidth}{lX}
\toprule
Experiment mode & Interpretation \\
\midrule
Inline literal input & 测量从 script text 解析 point lists、创建临时数据结构、分发 command 并执行 backend 的完整成本。该模式方便，但大输入时容易被 script parsing 主导。 \\
Prebuilt UVec without readback & 测量 handle lookup、device synchronization、output allocation 和 backend execution，同时将数据保持在 UVec form 中。该模式最接近 high-throughput script pipeline。 \\
Prebuilt UVec with readback & 额外加入将 output UVec 转为 script list 的成本。对于 distance matrices 这类 large outputs，readback 会主导总时间并掩盖 CUDA computation 的收益。 \\
\bottomrule
\end{tabularx}
\end{table}

消融实验给出了一条实践原则：FTCL scripts 应把 UVec handles 当作 long-lived data objects 使用，避免把 large intermediate results 转回 strings。这与 GPU programming 中的 data residency 原则一致。加速效果最强的情况是：数据只复制到 device 一次，随后由多个 commands 处理，最后再 read back。

这个解释也说明了为什么 CUDA 并不总是在小输入上更快。即使 kernel 本身很高效，整个 command path 仍然要支付 launch、synchronization、output allocation 和 possible readback 的 fixed costs。因此，break-even point 是 system-level property，而不仅是 algorithmic property。

Concurrent geometry experiment 如图 \ref{fig:concurrency-throughput} 和图 \ref{fig:concurrency-latency} 所示。CPU throughput 随 worker 数量增加而提升，直到 contention 出现。CUDA throughput 更平缓，因为多个 script workers 会竞争同一个 GPU 和 synchronization point。

\begin{figure}[H]
  \centering
  \includegraphics[width=0.78\textwidth]{geometry-concurrency-throughput.png}
  \caption{并发几何 workload 的吞吐量（Throughput of concurrent geometry workloads）}
  \label{fig:concurrency-throughput}
\end{figure}

直接结论：增加 script workers 只能在 shared runtime 或 GPU resource 成为 bottleneck 之前提升 throughput。

\begin{figure}[H]
  \centering
  \includegraphics[width=0.78\textwidth]{geometry-concurrency-latency.png}
  \caption{并发几何 workload 的 P95 延迟（P95 latency）}
  \label{fig:concurrency-latency}
\end{figure}

直接结论：contention 会推高 P95 latency，因此后续需要 dedicated scheduling、memory pooling 和 stream-based CUDA execution。

\section{小结}
本章从 correctness、performance 和 scalability 三个维度评估 FTCL。Semantic regression 和 CPU/GPU equivalence tests 证明系统具有可验证的正确性基础；single-GPU 和 multi-GPU experiments 显示 CUDA backend 在大量 independent work items 上具有明显优势；ablation 和 concurrency experiments 则指出真正的瓶颈经常来自 allocation、readback、locking 和 scheduling。

The evaluation is deliberately end-to-end. The main lesson is that data residency, output size, and runtime synchronization dominate script-level heterogeneous performance.

\chapter{Applications and Discussion}

\section{Ocean-Oriented Spatial Computing Demo}
Figure \ref{fig:real-demo} presents an ocean-oriented spatial computing demo. Instead of using a game scene, the demo models a coastal application scenario that is closer to marine ranching, smart ports, ship navigation, and unmanned inspection. Sensor buoys, unmanned aerial vehicles, floating objects, ships, aquaculture cages, port obstacles, and sonar/radar coverage regions are all represented as geometric data.

In this scenario, points represent sensor locations, floating objects, target observations, and UAV positions. Circles represent radar, sonar, or camera coverage. Axis-aligned boxes represent ships, cages, restricted port regions, or rectangular obstacles. Batch distance matrices represent large-scale relationships between sensors and abnormal targets. With UVec, these objects can stay on CUDA devices across several commands instead of being repeatedly converted back to script strings.

\begin{figure}[H]
  \centering
  \includegraphics[width=0.90\textwidth]{geometry-real-demo.png}
  \caption{Ocean-oriented spatial computing scene controlled by FTCL scripts}
  \label{fig:real-demo}
\end{figure}

The script pipeline in the figure uses the same geometry commands as the benchmark section, but the meaning of the data changes. The command \texttt{geom nearest\_point} can find the closest abnormal floating object to a UAV or patrol vessel. The command \texttt{geom range\_count\_circle} can count how many targets fall inside a radar, sonar, or camera coverage region. The command \texttt{geom collision\_aabb} can check whether a ship bounding box overlaps a reef, cage, pier, or restricted region. The command \texttt{geom batch\_distance\_matrix} can compute large sensor-target distance relationships for later assignment or alert scoring.

\section{Marine Ranch Spatial Perception}
In a marine ranch, buoys and underwater sensors continuously observe water quality, fish cage status, floating debris, and abnormal surface targets. A simple FTCL script can load sensor locations and object detections into UVec handles, run range queries for each sonar coverage circle, and report the nearest abnormal object to each inspection point.

This scenario is suitable for FTCL because the input naturally forms large point sets. Many queries are independent: each sensor coverage region can be processed separately, and each target-to-sensor distance can be computed in parallel. CUDA is useful when the number of detections or sensors becomes large, while the script layer remains convenient for combining several operations into one pipeline.

\section{Smart Port Spatial Computing}
A smart port contains berths, cranes, containers, ship boxes, navigation channels, and restricted regions. These objects can be represented using points, line segments, polygons, and bounding boxes. FTCL can use \texttt{geom collision\_aabb} to check coarse collisions between ships and port obstacles, \texttt{geom range\_count\_rect} to count objects inside management regions, and \texttt{geom nearest\_point} to find the closest available resource or warning target.

The value here is not that FTCL replaces a complete port information system. Instead, it provides a lightweight spatial-computing layer. A user can prototype port-side spatial rules in a script, then move large batch operations to the CUDA backend without changing the high-level command structure.

\section{UAV Inspection of Floating Objects}
Unmanned aerial vehicles are often used to inspect nearshore water surfaces. The UAV trajectory can be stored as a point sequence, detected floating objects can be stored as another point set, and camera coverage can be represented as circle queries. FTCL can compute nearest abnormal targets, count detections in each frame coverage region, and build a distance matrix between UAV sample points and suspicious objects.

This is a good example of why script-level heterogeneous computing is useful. The high-level inspection logic changes frequently: thresholds, target classes, and coverage radii may change from one task to another. The geometry kernels, however, remain stable and can be accelerated on GPU.

\section{Ship Path Planning and Collision Checking}
For ship navigation, a planned path must be checked against static and dynamic obstacles. In the current FTCL geometry subset, coarse collision checking can be performed using AABB commands, while segment intersection and distance commands can support later path-level checks. A ship route can also be evaluated against aquaculture cages, reefs, port boundaries, and moving obstacle boxes.

This scenario points to a natural extension direction. After the current batch geometry commands, future work can add spatial grids, polyline distance, route risk scoring, and dynamic obstacle prediction. These extensions would make FTCL more directly useful for marine and port spatial computing tasks.

\section{Algorithm Visualization}
Figure \ref{fig:geom-visual} visualizes representative algorithms: point-in-polygon classification, nearest point, circle range query, and AABB collision. Such figures are useful in a thesis defense because they explain the workload without requiring the reader to inspect code.

\begin{figure}[H]
  \centering
  \includegraphics[width=0.88\textwidth]{geometry_algorithm_visualization.png}
  \caption{Visualization of representative geometry algorithms}
  \label{fig:geom-visual}
\end{figure}

\section{Difference from a Normal Geometry Library}
A normal geometry library exposes C++ functions such as \texttt{distance(a,b)}.
Another common example is \texttt{segments\_intersect(a,b,c,d)}.
FTCL is different in four ways.

\begin{itemize}
  \item It exposes geometry operations through scripts, so users can combine tasks without recompiling C++.
  \item It uses UVec handles to avoid repeatedly converting large arrays into script strings.
  \item It supports both CPU and CUDA backends under one command interface.
  \item It includes tests, random differential checks, equivalence checks, and generated figures for thesis-level evaluation.
\end{itemize}

Therefore, the project is not simply a geometry function collection. Its value is the integration of interpreter design, data management, heterogeneous computing, and experimental validation.

From an application perspective, this also changes the positioning of the project. FTCL is not only a computational geometry scripting system. It can be described as a prototype heterogeneous computing substrate for ocean and port spatial analysis. The geometry commands are the first algorithmic layer; UVec and CUDA provide the execution substrate; the script language provides the control layer.

\section{Limitations}
The current system has several limitations.

First, UVec is not yet fully aligned with \texttt{std::vector}. It supports cross-device storage and script handles, but it should add more familiar container operations such as iterators, capacity management, insertion, erasure, and stronger exception-safety guarantees.

Second, the CUDA backend is still simple. It does not yet use memory pools, asynchronous streams, advanced reductions, shared-memory tiling, or spatial indexes. These techniques are important for reducing overhead and improving performance.

Third, the token-stream parser is architecturally cleaner but not always faster than the legacy parser. Future work should reduce token copying, compress whitespace handling, avoid repeated lexing for sub-scripts, and introduce AST or bytecode execution.

Fourth, the benchmark figures should be repeated under a fixed hardware environment with more rounds and confidence intervals. Some CUDA break-even points can fluctuate because end-to-end timing includes allocation, synchronization, and system scheduling noise.

\section{Practical Value for the Graduation Project}
This project is suitable as a graduation thesis because it has both engineering depth and algorithmic content. The interpreter demonstrates systems programming. The geometry library demonstrates algorithm design. UVec and CUDA demonstrate heterogeneous computing. The benchmark pipeline demonstrates reproducibility. Together, these parts form a richer result than a single algorithm demo.


\backmatter
\JOUBackmatterChapter{Conclusion and Future Work}
This thesis completed an end-to-end system that starts from a Tcl-style interpreter and grows into a CPU/CUDA hybrid computational geometry scripting platform. FTCL now supports a practical language subset, thread channels, UVec cross-device vectors, CPU geometry algorithms, CUDA kernels, tests, benchmarks, and thesis-ready figures. The experiments show that the GPU backend is most useful when work items are numerous and independent, while small inputs and immediate readback may erase the benefit of parallel execution.

Future work should proceed in four directions. First, UVec should be further aligned with the standard vector interface, including iterators, capacity management, insertion, erasure, and stronger exception-safety guarantees. Second, the CUDA backend should add memory pools, asynchronous streams, and better batching to reduce allocation and synchronization costs. Third, the parser path should continue moving toward token-stream parsing, AST representation, bytecode execution, and expression-level Pratt parsing. Fourth, the geometry library can be extended with spatial indexes, dynamic nearest-neighbor search, polygon boolean operations, and more realistic application demos.

\JOUAcknowledgementChapter
感谢党和国家对我的关怀，感谢人民和社会对我的培养与支持。

感谢我的指导老师纪兆辉老师在论文写作过程中给予的悉心指导。

感谢汪前进老师、朱敏老师、胡文彬老师、李慧老师、贾冬宝老师、吴加莹老师、韩璧如老师等各位老师在本科阶段对我的教导与帮助。

感谢冷严伟、刘好、孙京权等一路同行的伙伴。愿我们在未来继续坚定理想、脚踏实地，成就更有意义的事业。感谢班长吕勇同学在班级事务中的辛勤付出与帮助。

感谢我的家人长期以来的理解、支持与陪伴。

周恩来总理少年时曾立下“为中华之崛起而读书”的志向。今后，我也将以此自勉，继续锤炼专业本领，坚守初心，努力以所学所长服务国家发展、回报人民和社会。

\nocite{*}
\JOUPrintBibliography

\appendix
\chapter{构建、测试与复现命令}

本项目已开源，完整代码、测试用例与图表生成脚本见 GitHub 仓库：
\begin{center}
  \url{https://github.com/homer-fang/ftcl}
\end{center}

克隆仓库后，可按下列命令完成构建、测试与实验复现。以下路径以 WSL 下的 \texttt{/mnt/d/ftcl/ftcl} 为例；若从 GitHub 克隆，请将工作目录替换为本地克隆路径。

\section{CPU 构建与测试}
\begin{lstlisting}[language=bash]
git clone https://github.com/homer-fang/ftcl.git
cd ftcl
cmake -S . -B build -DCMAKE_BUILD_TYPE=Release \
  -DBUILD_TESTS=ON \
  -DFTCL_ENABLE_CUDA=OFF \
  -DFTCL_CXX_STANDARD=20
cmake --build build -j8
ctest --test-dir build --output-on-failure
\end{lstlisting}

\section{CUDA 构建}
\begin{lstlisting}[language=bash]
cmake -S . -B /tmp/ftcl-build-geometry-cuda \
  -DCMAKE_BUILD_TYPE=Release \
  -DBUILD_TESTS=ON \
  -DFTCL_ENABLE_CUDA=ON \
  -DFTCL_CXX_STANDARD=20
cmake --build /tmp/ftcl-build-geometry-cuda -j8
\end{lstlisting}

\section{核心基准测试}
\begin{lstlisting}[language=bash]
./build/test/bench_ftcl ./docs/data/benchmarks
python3 ./docs/scripts/benchmarks/plot_benchmarks.py \
  ./docs/data/benchmarks ./docs/figures/benchmarks
\end{lstlisting}

\section{几何图表生成}
\begin{lstlisting}[language=bash]
python3 docs/scripts/geometry/generate_geometry_figures.py \
  --build-dir /tmp/ftcl-build-geometry-cuda

python3 docs/scripts/geometry/generate_geometry_study_addons.py \
  --build-dir /tmp/ftcl-build-geometry-cuda
\end{lstlisting}

\section{论文编译}
\begin{lstlisting}[language=bash]
latexmk -xelatex -interaction=nonstopmode -output-directory=build main.tex
\end{lstlisting}

按当前 \texttt{latexmk} 配置，生成的 PDF 输出为 \texttt{build/main.pdf}。


\end{document}
